\section{Description}
\label{sec:architecture}

%# Descrizione
%> Copiando da un paper di Cucinotta, qui descriverei solo il modello di cluster 
%    e l'API di ogni nodo sia per i client che per nodi.
%
%**No coordinazione.**
%Altre varie assunzioni:
%- Modello di cluster (possono esplodere per tempo finito).
%- Modello di banda.
%- No Consistenza: si fa con un layer aggiuntivo
%- API client + tabelle senza nome

\Kitsurai stores users' data in a set of \emph{tables}, each identified by a
UUID~\cite{rfc9562}.
Every table contains key-value pairs, where both keys and values are
uninterpreted arrays of bytes.

A \Kitsurai system operates on a \emph{cluster} composed of a static set of
nodes, a \emph{node} corresponds to an instance of the program running on a
physical or virtual machine.
Each node may be temporarily unavailable due to failures or network partitions,
the system is designed to be resilient to such events and to continue operating
as long as a sufficient number of nodes are available.
Each node can be assigned to an \emph{Availability Zone} (see
Section~\ref{subsec:availability-zones}).

\subsection{Tables}
\label{subsec:tables}

Each table has four parameters specified by the user at creation time:
\begin{itemize}
    \item \texttt{bandwidth} or $B$: the number of requests, as
        perceived by the client, that the system must guarantee for the table,
        expressed in requests per second.
    \item $N$: the replication factor for each key-value pair.
    \item \texttt{read\_concern} or $R$: the minimum number of nodes
        that must successfully respond to a read request.
    \item \texttt{write\_concern} or $W$: the minimum number of nodes
        that must successfully respond to a write request.
\end{itemize}

Parameters must satisfy the constraints $R, W \in \{ 1, 2, ..., N \}$ and
$B > 0$.

Additionally, the system computes four parameters automatically:
\begin{itemize}
    \item \texttt{id}: a UUID identifying the table, generated at creation.
    \item \texttt{responsible\_nodes}: the list of nodes on which the table's
        key-value pairs are distributed, chosen to guarantee replication and
        data availability.
    \item $m$: the number of \texttt{responsible\_nodes}.
    \item $s_i$: the bandwidth allocated by node $i$ for the table.
    \item $S$: the total bandwidth allocated for the table, computed as $S =
        \sum s_i$ or equivalently $S = B \cdot N$.
\end{itemize}

\subsection{Consistency and Client-Visible Semantics}
\label{subsec:consistency}

Our store is a best-effort, partition-tolerant key–value service with explicit
client-tunable quorums. We make two deliberate design choices:

\begin{itemize}
    \item \textbf{No replica reconciliation.}
        Item payloads receive \emph{no} server-side conflict resolution and
        \emph{no} catch-up to the latest value for restarting or recovering
        replicas.
        Replicas that miss a write may indefinitely retain a divergent value
        until a later client write updates them.
    \item \textbf{Multi-value reads.}
        Reads may return multiple distinct versions of the same key.
        The store does not attempt to infer a single ``latest'' value, this is
        left to the application.
\end{itemize}

\paragraph{Quorum parameters.}
We use lowercase $N, R, W$ for the number of replicas, read quorum, and write
quorum.
A write to key $k$ is considered \emph{committed} once any $W$ replicas
acknowledge persisting the value; the system does not subsequently propagate
that value to the remaining $N - W$ replicas.
A read for key $k$ completes once responses from any $R$ replicas are received.

\paragraph{Quorum intersection.}
If $R + W > N$, then any set of $R$ replicas intersects any set of $W$
replicas.
Consequently, every read that gathers $R$ replies contains \emph{at least one}
payload written by any previously completed write.
However, in the presence of concurrent or failed writes, multiple distinct
payloads can appear in the read result, and the store offers no mechanism to
determine which is the ``latest'' value.

\subsection{Client API}
\label{subsec:interface}

The system exposes the public API through a REST interface over HTTP\@.
Since every node maintains a copy of the metadata for all tables, the client
can send requests to any node in the cluster, which acts as a \emph{router} to
forward the request to the appropriate nodes.

\subsubsection{\texttt{table-create}}
\label{subsubsec:table-create}

The client sends four parameters: $B$, $N$, $R$, and $W$.
The system allocates enough bandwidth on the nodes to satisfy the requested
parameters, creates the table, and returns the generated \texttt{id} of the new
table.
If it is not possible to satisfy the requested parameters, the operation fails
and returns an error.

\subsubsection{\texttt{table-delete}}
\label{subsubsec:table-delete}

The client sends the \texttt{id} of the table to be deleted.
The system marks the table as deleted and frees the allocated resources.

\subsubsection{\texttt{item-get}}
\label{subsubsec:item-get}

The client sends the \texttt{id} of the table and the \texttt{key} to be
retrieved.
The system forwards the request to the responsible nodes and waits for at least
$R$ successful responses within a system-wide configurable timeout.
If the operation succeeds, it returns a list of the values received from the
nodes; otherwise, it returns an error with the successfully collected values.

\subsubsection{\texttt{item-set}}
\label{subsubsec:item-set}

The client sends the \texttt{id} of the table, the \texttt{key}, and the
\texttt{value} to be set.
The system forwards the request to the responsible nodes and waits for at least
$W$ successful responses within a system-wide configurable timeout.
If the operation succeeds, it returns a success message; otherwise, it returns
an error.

\begin{figure}
    \includegraphics[width=\columnwidth]{03-description/architecture}
    \caption{\Kitsurai architecture overview. Clients can send requests to any
        node, which routes them to the responsible nodes for the requested
        table.}
    \label{fig:architecture}
\end{figure}

\subsection{Availability Zones}
\label{subsec:availability-zones}

To improve fault tolerance and ensure continuous data accessibility, \Kitsurai
adopts a replication strategy based on \emph{Availability Zones} (AZs).
Each node within the cluster is associated with an Availability Zone
identifier, a symbolic label representing either a physical or a logical
grouping of nodes, such as distinct data centers, server racks, or virtual
clusters distributed across geographical regions.

During table creation, the system enforces a zone-aware placement policy that
guarantees every key–value pair is replicated on $N$ nodes belonging to $N$
distinct Availability Zones.
Consequently, no two replicas of the same item reside within the same zone.

This design minimizes the likelihood of correlated failures leading to data
unavailability and ensures that a table remains operational even in the event
of a complete zone outage.
As long as the number of available zones satisfies the quorum parameters $(R,
W)$ defined for the table, \Kitsurai continues to serve read and write requests
using replicas from these zones.

\subsection{Throughput Throttling}
\label{subsec:throughput-throttling}

In a multi-tenant system such as \Kitsurai, it is necessary to prevent a single
client from monopolizing resources and compromising the performance of others.
For this reason, \Kitsurai implements a throughput throttling mechanism that
guarantees each table the bandwidth requested at table creation, even in the
presence of high and concurrent loads.

\subsection{Bandwidth Allocation Model}
\label{subsec:bandwidth-model}

For bandwidth allocation, \Kitsurai adopts a very simple model: it assumes that
the maximum amount of operations per second a machine can execute is fixed and
independent of other factors such as the size of the data being read or
written.

For this reason, each node only needs to know its maximum bandwidth, measured
beforehand using standard profiling operations, to determine whether it is able
to handle the load requested by a table.
For this purpose, SSD disks must be used, as they provide constant and
predictable performance, unlike HDDs, which are subject to unpredictable
slowdowns due to seek operations and access latency.
